% tab:sensitivity
\begin{table}[htbp]
\centering\footnotesize
\caption{Sensitivity of the Latvian series to the main sources of bias}
\label{tab:sensitivity}
\begin{threeparttable}
\begin{tabular}{p{4.6cm}lp{3.8cm}rr}
\toprule
Source of bias & Years & Sectors affected & Bias & Max.\ impact \\
& & & (\%) & on GDP (\%) \\
\midrule
Incorrect service-sector evaluation & 1919--40 & Services & $\pm$10 & $\pm$3 \\
Fixed value-added weights (single deflation) & 1919--40 & All & $\pm$2 & $\pm$2 \\
Labour productivity and real-wage estimates & 1919--29 & Industry & $\pm$10 & $\pm$2 \\
Missing annual product-tax data & 1919--40 & Product taxes & $\pm$9 & $\pm$1 \\
No productivity adjustment in employment series & 1919--40 & Public services & $\pm$10 & $\pm$2 \\
Population movement as proxy & 1919--40 & Various service subsectors & $\pm$30 & $\pm$1 \\
\midrule
Incorrect service-sector evaluation & 1940--90 & Services (excl. transport) & $\pm$10 & $\pm$4 \\
Missing product-tax data & 1940--90 & Total & $-$0.5 & $-$0.5 \\
Fixed value-added weights (single deflation) & 1940--90 & All & $-$14 & $-$14 \\
Hidden inflation in trade series & 1945--90 & Trade, finance and insurance & $+$59 & $+$2 \\
Unknown Soviet military size & 1941--89 & Defence & $\pm$20 & $\pm$7 \\
No productivity adjustment in employment series & 1940--90 & Housing, education, health, science, administration & $-$147 & $-$10 \\
\bottomrule
\end{tabular}
\begin{tablenotes}[para,flushleft]
\footnotesize
\textit{Notes:} `Bias' is the assumed maximum error in the affected indicator, taken as 10 per cent in either direction unless an independent estimate exists; `Max.\ impact' is that error weighted by the affected sectors' share of GDP in the year in which it is largest. A minus sign means that growth in the new series may be understated, so that backcast levels are too high, and vice versa. Fixed value-added weights: Swedish single- versus double-deflated accounts \parencites[19]{swedengdp_15602010}[148--53]{lobellschonkrantz2008}; hidden inflation: difference between official and estimated real inflation in Soviet retail trade \parencite[607]{hidden_inflation_sov}; product taxes: in 1980--90 the ratio of product taxes less subsidies to value-added was 0.5 per cent in Bulgaria and Ukraine, the only socialist economies with data \parencite{wdi2024}; service productivity: British service-sector labour productivity rose by 147 per cent between 1948 and 1990 \parencite{statisticsUKlabour}. \textit{Sources:} \textcite[149--56]{klimantas2025thesis}; this paper.
\end{tablenotes}
\end{threeparttable}
\end{table}
